\scanpage{185}

De plus, on a une structure \unsure{polarisée} \inserted{$\pi$} sur $G^0_{\mathbb R}$, qui est \deleted{\unsure{déf}} supposée définie par un hom
$\deleted{\unsure{\mu_{2\mathbb C}}}\quad i:\mu_{2\mathbb C}\unsure{-|\to}G^0_{\mathbb C}$
(comme on voit en envoyant un \unsure{tore} \inserted{point d'ordre 2} dedans). D'autre part, $G(s')_{\mathbb C}$ est muni de $i(s'):G_{m\mathbb C}\to G^{(s')}_{\mathbb C}$, et on obtient par conséquent
\diagram{\[
\begin{tikzcd}[column sep=large]
\unsure{v|v_q}(s')=u(s')_{\mathbb C}\circ(i(s')\mid\mu_2):\ \mu_{2\mathbb C} \arrow[r] & G^{\unsure{\xi}(s')}_{\mathbb C}
\end{tikzcd}
\]}

\deleted{D'autre part} Ceci posé, il faut que $\pi^{\unsure{\xi}(s')}$ (polarisation de $G^{\unsure{\xi}(s')}_{\mathbb R}$ déduite de $\pi$ sur $\mathbb C$ par transport) soit celle définie par les $\mu_{2\mathbb C}$.

Par conséquent, pour $s'$ \unsure{variable|variant} de façon que $\unsure{\xi}(s')$ soit constant, les polarisations \inserted{de $G^{\unsure{\xi}(s')}_{\mathbb R}$} associées aux différents $v(s')$ doivent être identiques -- ce qui donne une condition non tautologique sur la distribution des \deleted{\unsure{ces}} $u(s')$. [Noter cependant que l'ens des \deleted{classes de} polarisations d'un gpe réel \inserted{$H$}, définies en termes d'un hom $\mu_{\mathbb C}\to H_{\mathbb C}$, est \uline{fini}, [si $G$ est commutatif en tous cas !], donc la restriction n'est pas tellement draconienne.
\margin{? à vérifier}

Enfin, conformément aux conjectures de Tate \unsure{magnifiées}, La connaissance des $u(s')$ permet de retrouver la structure d'effectivité dans la catégorie

\margin{184}
